\uuid{3kdw}
\chapitre{Statistique}
\niveau{L2}
\module{Analyse de données}
\sousChapitre{Estimation}
\titre{Lois statistiques et étude d'estimateurs}
\theme{statistiques, estimateurs}
\auteur{Maxime NGUYEN}
\datecreate{2022-09-07}
\organisation{AMSCC}
\difficulte{}
\contenu{

\texte{$\mathcal N(\mu,\sigma)$ désigne la loi normale de moyenne $\mu$ et d’écart-type $\sigma$. Les quantiles sont à gauche : $q_{\nu;\gamma}$ pour la loi $\chi^2(\nu)$, avec une probabilité $\gamma$ à gauche du quantile.}
\texte{Les variables $X_1,\ldots,X_5$ sont indépendantes, de loi $\mathcal N(\mu,\sigma)$, $\sigma>0$. On pose
\begin{align*}
T_1&=\frac15\sum_{i=1}^5X_i, & T_2&=\frac{X_1+X_2}{5}+\frac{X_3+X_4+X_5}{4},\\
T_3&=\frac{2X_1+3X_2}{10}+\frac{X_3+2X_4+X_5}{8},\\
U&=\frac{\sum_{i=1}^5(X_i-\mu)^2}{\sigma^2},&V&=\frac{\sum_{i=1}^5(X_i-T_1)^2}{\sigma^2}.
\end{align*}}

	\begin{enumerate}
		\item \question{Quelle est la loi suivie par la variable $X_1-X_2$ ? Justifier.}
		\reponse{
			D'après le cours, $X_1-X_2$ suit une loi normale d'espérance $\mathbb{E}(X_1-X_2) = \mu - \mu = 0$. Par indépendance, $\operatorname{Var}(X_1-X_2) = \operatorname{Var}(X_1)+(-1)^2\operatorname{Var}(X_2) = 2\sigma^2$. Ainsi, $X_1-X_2\sim\mathcal N(0,\sigma\sqrt2)$.
		}
		\item \question{On cherche à estimer $\mu$ à l'aide des estimateurs $T_1$, $T_2$, $T_3$. \'Etudier leur biais et comparer l'efficacité des estimateurs sans biais.}
		\reponse{ Par linéarité de l'espérance, on calcule $\mathbb{E}(T_1) = \frac{5\mu}{5} = \mu$, $\mathbb{E}(T_2) = \frac{2\mu}{5}+\frac{3\mu}{4}=\frac{23\mu}{20}$, $\mathbb{E}(T_3) = \mu$. Par conséquent, $b(T_1)=b(T_3)=0$ et $b(T_2) = \mathbb{E}(T_2)-\mu = \frac{3\mu}{20}$.

			Pour comparer l'efficacité des deux estimateurs sans biais, on calcule leur EQM (ce qui revient à calculer leur variance.) Par indépendance des variables, on a :

			\[
\operatorname{Var}(T_1)=\frac{\sigma^2}{5}.
\]
\begin{align*}
\operatorname{Var}(T_3)
&=\left(\frac4{100}+\frac9{100}+\frac1{64}+\frac4{64}+\frac1{64}\right)\sigma^2\\
&=\frac{179\sigma^2}{800}.
\end{align*}
Comme $\frac15=\frac{160}{800}<\frac{179}{800}$, le plus efficace est donc l'estimateur $T_1$ qui est la moyenne empirique.

		}
		\item \question{Quelle est la loi suivie par la variable $U$ ? la variable $V$ ? justifier.}
		\reponse{$U = \sum_{i=1}^{5}  \left(\frac{X_i-\mu}{\sigma}\right)^2$ ; or les $X_i$ sont des variables aléatoires indépendantes et $\frac{X_i-\mu}{\sigma}$ suit une loi $\mathcal{N}(0,1)$ donc par définition, $U$ suit une loi de $\chi^2(5)$.

			De plus, $T_1 = \overline{X}$ est l'estimateur de moyenne empirique donc d'après le théorème de Fisher, $V$ suit une loi de $\chi^2(5-1) = \chi^2(4)$ et est indépendant de $T_1$. }
		\item \question{Déterminer $x \in \R$ tel que $\PP(U>x) = 0{,}05$.}
		\reponse{Comme $U$ suit une loi continue, $\PP(U>x)=0{,}05$ équivaut à $\PP(U\leq x)=0{,}95$. On lit dans la table de loi le quantile $x=q_{5;0{,}95}\simeq11{,}07$. }
		\item \begin{samepage}
\question{En utilisant $T_1$ et $U$, construire une variable $Y$ qui suive une loi de Student dont on précisera le paramètre.}
		\indication{Décomposer la somme des carrés autour de $\mu$ en utilisant la moyenne empirique, puis vérifier l'indépendance nécessaire à la construction de Student.}
\reponse{\begin{minipage}{\linewidth}
On pose
\[
Z=\frac{T_1-\mu}{\frac{\sigma}{\sqrt5}}
=\frac{\sqrt5(T_1-\mu)}{\sigma}.
\]
La moyenne empirique $T_1$ suit la loi $\mathcal N(\mu,\frac{\sigma}{\sqrt5})$, donc $Z\sim\mathcal N(0,1)$.

Pour conserver une construction à partir de $T_1$ et de $U$, on décompose la somme des carrés. Puisque $\sum_{i=1}^5(X_i-T_1)=0$, on a
\begin{align*}
\sum_{i=1}^5(X_i-\mu)^2
&=\sum_{i=1}^5\bigl((X_i-T_1)+(T_1-\mu)\bigr)^2\\
&=\sum_{i=1}^5(X_i-T_1)^2+5(T_1-\mu)^2.
\end{align*}
Après division par $\sigma^2$, on obtient $U=V+Z^2$, donc $U-Z^2=V$.

Le théorème de Fisher donne $V\sim\chi^2(4)$ et l'indépendance entre $V$ et $T_1$, donc entre $V$ et $Z$. Ainsi,
\[
Y=\frac{Z}{\sqrt{\frac{U-Z^2}{4}}}
=\frac{Z}{\sqrt{\frac V4}}
\sim St(4).
\]
La variable construite suit une loi de Student à 4 degrés de liberté.
La construction $\frac{Z}{\sqrt{\frac U5}}$ ne donne pas une loi de Student à 5 degrés de liberté, car $U=V+Z^2$ dépend de $Z$.
\end{minipage}}
\end{samepage}
	\end{enumerate}

	}
